\section{Einleitung}

In der folgenden Sektion werden wir nun das Wissen aus den vorherigen Kapiteln nutzen und
einmal selbst eine Faltung zweier Audiosignale durchführen.
Dafür verwenden wir die zuvor vorgestellte Gleichung der diskreten Faltung \ref{scaled_unit_sum}
aus dem Kapitel \ref{sec:convolution}:

\begin{equation}
    y[n] = \sum_{k=0}^{M-1}h[k] \cdot x[n-k]
    \label{eq:directConv}
\end{equation}
Wobei:
\begin{itemize}
    \item[$n$]: Folgenindex (Maximum ist Länge des Ausgangssignals)
    \item[$M$]: Länge der Impulsantwort 
\end{itemize}

Für die konkrete Umsetzung als Programm nutzen wir die Sprache \textit{C++}.
Es sollen zuerst die diskreten Abtastwerte der jeweiligen Audiodateien in Arrays
geladen werden. Danach soll die Faltung durchgeführt und anschließend das Ergebnis als neue
Audiodatei im Hintergrundspeicher abgelegt werden.

Für das Laden und speichern von Audiodateien im .wav-Format nutzen wir ein bekanntes Framework
namens \textit{JUCE} \cite{juce}.
\textit{JUCE} bietet noch sehr viele weitere Funktionalitäten, wie bspw. Klassen
zum Entwurf von Benutzeroberflächen oder Schnittstellen zur Entwicklung von \textit{Plug-ins},
die direkt in Programmen wie \textit{Cubase} eingebunden werden können.
Die Implementierung solcher Anwendungen würden hierbei jedoch absolut den Rahmen sprengen, weshalb
wir uns hier lediglich auf die Faltung an sich beschränken.
Bei Interesse sei also hiermit auf die offizielle Website und Dokumentation von \textit{JUCE} verwiesen: \cite{juce}.
